\changecolours{ScoutGreen}
\section[Conclusão e Próximos Passos]{Conclusão, Atividades e Conexão com a Aula 06}

% ------------------------------------------------------------------------------
\twocolframe[title={Síntese Conceitual da Aula 05},leftcol=7.1cm,rightcol=6.9cm]{%
  \alert{Checklist de Competências Consolidadas}\par
  \itemseps[topsepi=1mm,itemsepi=1.5mm,labelsep=2mm,leftmargini=3.5mm]
  \begin{itemize}
    \item \textbf{Definição Formal:} Relacionamentos como subconjunto do produto cartesiano ($R \subseteq E_1 \times E_2 \times \dots$).
    \item \textbf{Grau e Papéis:} Diferenciação clara de unário, binário e ternário, com papéis explícitos para desambiguação.
    \item \textbf{Cardinalidade:} Domínio absoluto das razões $1:1$, $1:N$ e $N:M$.
    \item \textbf{Participação:} Linha dupla para dependência de existência (total) e simples para opcional (parcial).
    \item \textbf{Atributos Descritivos:} Alocação no losango quando a informação pertence ao fato da ligação.
  \end{itemize}
}{%
  \alert{Boas Práticas de Modelagem no MER}\par
  \itemseps[topsepi=1mm,itemsepi=1.5mm,labelsep=2mm,leftmargini=3.5mm]
  \begin{itemize}
    \item \textbf{Nomenclatura Significativa:}
      Utilize verbos expressivos no infinitivo ou presente que comuniquem a regra de negócio com precisão.
    \item \textbf{Eliminação de Redundâncias:}
      Não crie relacionamentos diretos que possam ser deduzidos com segurança através de outros vínculos existentes.
    \item \textbf{Rotulação Semântica de Papéis:}
      Sempre adicione papéis em relacionamentos unários e ternários para evitar ambiguidades de interpretação.
    \item \textbf{Validação Prévia com o Negócio:}
      Confirme com o usuário final cada restrição de cardinalidade antes de iniciar o esquema relacional.
  \end{itemize}
}

% ------------------------------------------------------------------------------
\twocolframe[title={Ponte Pedagógica: O Que Vem na Aula 06?},leftcol=7cm,rightcol=7cm]{%
  \alert{O Próximo Encontro: Aula 06 (14/10/2026)}\par
  \itemseps[topsepi=1mm,itemsepi=2mm,labelsep=2mm,leftmargini=3.5mm]
  \begin{itemize}
    \item \textbf{Entidades Fracas (\textit{Weak Entities}):} O que fazer quando uma entidade não possui atributos suficientes para ter chave própria?
    \item \textbf{Chave Parcial / Discriminador:} Identificação subordinada a uma entidade proprietária (sublinhado tracejado).
    \item \textbf{Relacionamento Identificador:} O losango duplo no DER de Peter Chen.
    \item \textbf{Modelo EER (Estendido):} Introdução a hierarquias de herança, especialização e generalização de dados.
  \end{itemize}
}{%
  \alert{Atividades de Fixação e Projeto AVP1}\par
  \itemseps[topsepi=1mm,itemsepi=2mm,labelsep=2mm,leftmargini=3.5mm]
  \begin{itemize}
    \item \textbf{Exercícios da Aula 05:} Resolver as 5 questões conceituais do plano de aula no caderno de estudos.
    \item \textbf{Execução do Notebook:} Executar o script Python didático para observar as regras e mensagens de bloqueio.
    \item \textbf{Alinhamento da AVP1:} Esboçar as primeiras entidades e relacionamentos do minimundo do seu projeto semestral.
  \end{itemize}
}
